%% LyX 2.4.4 created this file.  For more info, see https://www.lyx.org/.
%% Do not edit unless you really know what you are doing.
\documentclass[english]{article}
\usepackage[T1]{fontenc}
\usepackage[latin9]{inputenc}
\usepackage{units}
\usepackage{enumitem}
\usepackage{amsmath}
\usepackage{amsthm}
\usepackage{amssymb}
\usepackage{geometry}
\geometry{verbose,tmargin=1in,bmargin=1in,lmargin=1in,rmargin=1in}
\PassOptionsToPackage{normalem}{ulem}
\usepackage{ulem}

\makeatletter
%%%%%%%%%%%%%%%%%%%%%%%%%%%%%% Textclass specific LaTeX commands.
\numberwithin{equation}{section}
\numberwithin{figure}{section}
\newlength{\lyxlabelwidth}      % auxiliary length 

%%%%%%%%%%%%%%%%%%%%%%%%%%%%%% User specified LaTeX commands.
\usepackage{fancyhdr}
\usepackage{lastpage}
\pagestyle{fancy}
\usepackage{enumitem}
\usepackage{ifthen}
\everymath=\expandafter{\the\everymath\displaystyle}
%\DeclareMathSizes{10}{10}{10}{10}
\usepackage{multicol}

\usepackage{tikz}
\usepackage{tikz-3dplot}
\usetikzlibrary{calc,arrows}

\usetikzlibrary{patterns}
\usetikzlibrary{plotmarks}
\usepackage{pgfplots}
\pgfplotsset{compat=newest}

\usepackage{calc}
\usepackage[nomessages]{fp}

\usepackage{datetime}
% Added by lyx2lyx
\setlength{\parskip}{\smallskipamount}
\setlength{\parindent}{0pt}

\makeatother

\usepackage{babel}
\begin{document}
\renewcommand{\author}{Dr.\,James Ashe}

\newcommand{\classNum}{336}

\newcommand{\classSec}{A}

\newcommand{\examType}{Exam 3}

\newcommand{\nameType}{Solutions}

\renewcommand{\date}{\formatdate{20}{11}{2013}}

%%First page's header%%%%%%%%%%%%%%%%%%%%%%
\fancypagestyle{firststyle} %%header settings for first pagr
{
	\fancyhead[L]{MTH \classNum{}(\classSec{}) \author{}\\
	\textbf{\examType{}} ~\date{}}
	\fancyhead[C]{\textbf{\nameType}}
	\fancyhead[R]{}
	\setlength{\headsep}{14pt}
}
%%%%%%%%%%%%%%%%%%%%%%%%%%%%%%%%%%%%%%%%%%

%%Next page's header%%%%%%%%%%%%%%%%%%%%%%%
\setlist{nolistsep}
\lhead{MTH \classNum{}(\classSec{})} 
\chead{\textbf{\nameType}} 
\cfoot{\thepage{}~of~\pageref{LastPage}} 
\setlength{\headsep}{5pt} 
%%%%%%%%%%%%%%%%%%%%%%%%%%%%%%%%%%%%%%%%%%%

%%Various settings; see the preamble for other settings%%%%%
%\setenumerate[0]{leftmargin=12pt,itemindent=0pt,labelwidth=0pt}
%\setlist{topsep=.5pt}
\renewcommand{\labelenumi}{\textbf{\arabic{enumi}.}}
\renewcommand{\labelenumii}{\textbf{\alph{enumii}.}}
\thispagestyle{firststyle}
%%%%%%%%%%%%%%%%%%%%%%%%%%%%%%%%%%%%%%%%%%

\newcommand{\xmax}{0}
\newcommand{\xmin}{-\xmax}
\newcommand{\xstep}{0}
\newcommand{\ymax}{0}
\newcommand{\ymin}{-\ymax}
\newcommand{\ystep}{0} 
\newcommand{\dspace}{\vspace{1cm}}
\newcommand{\xa}{0}
\newcommand{\xb}{0}
\newcommand{\ya}{1}
\newcommand{\yb}{1}

\newcommand{\xspace}{\vspace{0cm}}
\newcommand{\Xspace}{\vfill}

\newcommand{\vectors}[2]{\textbf{#1}_{1},\textbf{#1}_{2},\dots,\textbf{#1}_{#2}}
\newcommand{\vectorsN}[1]{\textbf{#1}_{1},\textbf{#1}_{2},\dots,\textbf{#1}_{n}}
\newcommand{\R}[1]{\mathbb{R}^{#1}}

\newcounter{probs}

\small

\textbf{Work must be shown to receive credit. Each problem is worth
5 pts unless stated otherwise.}
\begin{enumerate}[resume]
\item Use properties of determinants to calculate the determinant of the
following matrices:\begin{multicols}{3}

\begin{enumerate}
\item $\begin{bmatrix}\begin{array}{rrrr}
3 & 0 & 0 & 0\\
4 & 3 & 0 & 0\\
5 & 2 & 1 & 0\\
9 & 8 & 9 & 1
\end{array}\end{bmatrix}$
\begin{eqnarray*}
\det A & = & (3)(3)(1)(1)\\
 & = & 9
\end{eqnarray*}
\xspace \columnbreak\addtocounter{probs}{1}
\end{enumerate}
\begin{enumerate}[resume]
\item $\begin{bmatrix}\begin{array}{rrrr}
1 & 1 & 1 & 1\\
0 & 0 & 2 & 3\\
0 & 2 & 2 & 2\\
0 & 0 & 0 & 4
\end{array}\end{bmatrix}$
\begin{eqnarray*}
\det A & = & -\det\begin{bmatrix}\begin{array}{rrrr}
1 & 1 & 1 & 1\\
0 & 2 & 2 & 2\\
0 & 0 & 2 & 3\\
0 & 0 & 0 & 4
\end{array}\end{bmatrix}\\
 & = & (1)(2)(2)(4)\\
 & = & 16
\end{eqnarray*}
\xspace \columnbreak\addtocounter{probs}{1}
\item $\begin{bmatrix}\begin{array}{rrrr}
8 & 3 & 5 & 6\\
3 & 9 & \nicefrac{1}{2} & 7\\
0 & 0 & 0 & 0\\
2 & 4 & 6 & 7
\end{array}\end{bmatrix}$
\begin{eqnarray*}
\det A & = & 0
\end{eqnarray*}
\xspace \columnbreak\addtocounter{probs}{1}
\end{enumerate}
\end{multicols}
\end{enumerate}
\vspace{-1cm}
\begin{enumerate}[resume]
\item Let $A=\begin{bmatrix}\begin{array}{rrr}
1 & 2 & 3\\
0 & 4 & 5\\
1 & 2 & 3
\end{array}\end{bmatrix}$ 

\begin{enumerate}
\item Calculate the determinant of $A$ using cofactors. 

\begin{enumerate}
\item [](using row 2)$\det A=-0+4\det\begin{bmatrix}\begin{array}{rr}
1 & 3\\
1 & 3
\end{array}\end{bmatrix}-5\det\begin{bmatrix}\begin{array}{rr}
1 & 2\\
1 & 2
\end{array}\end{bmatrix}=0+0+0=0$\addtocounter{probs}{1}
\end{enumerate}
\end{enumerate}
\begin{enumerate}[resume]
\item Find $C_{32}$, i.e., $A$'s cofactor for the $3^{\mbox{rd}}$ row
and $2^{\mbox{nd}}$ column.

\begin{enumerate}[resume]
\item []$C_{32}=(-1)^{(3+2)}\det\begin{bmatrix}\begin{array}{rr}
1 & 3\\
0 & 5
\end{array}\end{bmatrix}=-5$\addtocounter{probs}{1}
\end{enumerate}
\item Does ${\displaystyle \sum_{j=1}^{3}a_{1j}C_{1j}=\sum_{j=1}^{3}a_{3j}C_{3j}}$
for any $3\times3$ matrix?

\begin{enumerate}[resume]
\item []Yes. The determinant is the same regardless of which row is used.\addtocounter{probs}{1}
\end{enumerate}
\end{enumerate}
\end{enumerate}
\vfill
\begin{enumerate}[resume]
\item Let $\begin{bmatrix}\begin{array}{rr}
2 & 3\\
4 & 7
\end{array}\end{bmatrix}\begin{bmatrix}\begin{array}{r}
x_{1}\\
x_{2}
\end{array}\end{bmatrix}=\begin{bmatrix}\begin{array}{r}
1\\
1
\end{array}\end{bmatrix}$.\begin{multicols}{2}

\begin{enumerate}
\item Find $x_{1}$ using Cramer's rule.

\begin{enumerate}
\item []$\det\begin{bmatrix}\begin{array}{rr}
2 & 3\\
4 & 7
\end{array}\end{bmatrix}=2$ and $\det\begin{bmatrix}\begin{array}{rr}
1 & 3\\
1 & 7
\end{array}\end{bmatrix}=4$,\\
so $x_{1}=\frac{4}{2}=2$\Xspace\addtocounter{probs}{1}
\end{enumerate}
\end{enumerate}
\begin{enumerate}[resume]
\item Find $x_{2}$ using Cramer's rule.

\begin{enumerate}
\item []$\det\begin{bmatrix}\begin{array}{rr}
2 & 3\\
4 & 7
\end{array}\end{bmatrix}=2$ and $\det\begin{bmatrix}\begin{array}{rr}
2 & 1\\
4 & 1
\end{array}\end{bmatrix}=-2$,\\
so $x_{1}=\frac{-2}{2}=-1$\Xspace\addtocounter{probs}{1}
\end{enumerate}
\end{enumerate}
\end{multicols}\vfill
\item Suppose that $A$ and $B$ are $n\times n$ matrices such that $\det A=5$
and $\det B=-2$.

\begin{enumerate}
\item Find $\det AB$.

\begin{enumerate}
\item []$\det AB=\det A\det B=(5)(-2)=-10$\xspace\addtocounter{probs}{1}
\end{enumerate}
\end{enumerate}
\begin{enumerate}[resume]
\item Prove that $\left(AB\right)^{-1}$ exists. 

\begin{enumerate}[resume]
\item []From above,$\det AB=-10\neq0$. So $AB$ is invertible.\xspace\addtocounter{probs}{1}
\end{enumerate}
\item Find $\det\left(AB\right)^{-1}$.

\begin{enumerate}[resume]
\item $\det(AB)^{-1}=\frac{1}{\det(AB)}=-\frac{1}{10}$\xspace\addtocounter{probs}{1}\vfill
\end{enumerate}
\end{enumerate}
\item Find the characteristic polynomials of the following matrices.

\begin{enumerate}
\item $\begin{bmatrix}\begin{array}{rrrr}
4 & 4 & 1 & 1\\
0 & 3 & 2 & 1\\
0 & 0 & 2 & 3\\
0 & 0 & 0 & 1
\end{array}\end{bmatrix}$

\begin{enumerate}
\item []$p_{A}(t)=\det(A-tI)=\det\begin{bmatrix}\begin{array}{rrrr}
4-t & 4 & 1 & 1\\
0 & 3-t & 2 & 1\\
0 & 0 & 2-t & 3\\
0 & 0 & 0 & 1-t
\end{array}\end{bmatrix}=(4-t)(3-t)(2-t)(1-2)$\addtocounter{probs}{1}\newpage
\end{enumerate}
\end{enumerate}
\begin{enumerate}[resume]
\item $\begin{bmatrix}\begin{array}{rrr}
-1 & -2 & 1\\
1 & 3 & 2\\
0 & 2 & 1
\end{array}\end{bmatrix}$

\begin{enumerate}[resume]
\item []Using cofactors, $p_{A}(t)=\det(A-tI)=\det\begin{bmatrix}\begin{array}{rrr}
-1-t & -2 & 1\\
1 & 3-t & 2\\
0 & 2 & 1-t
\end{array}\end{bmatrix}=0-2\det\begin{bmatrix}\begin{array}{rr}
-(1+t) & 1\\
1 & 2
\end{array}\end{bmatrix}+(1-t)\det\begin{bmatrix}\begin{array}{rr}
-(1+t) & -2\\
1 & (3-t)
\end{array}\end{bmatrix}=-2\left[-2(1+t)-1\right]+(1-t)\left[-(1+t)(3-t)+2\right]$ which simplifies to $p_{A}(t)=-t^{3}+3t^{2}+3t+5$.\Xspace\addtocounter{probs}{1}
\end{enumerate}
\end{enumerate}
\item Find the (real) eigenvalues of the following matrices.

\begin{enumerate}
\item Let $A$ be a $3\times3$ matrix and $p_{A}(t)=(t+1)(t-2)(t+3)$.

\begin{enumerate}
\item []Eigenvalues: $-1,\,2,\,-3$\Xspace\addtocounter{probs}{1}
\end{enumerate}
\item Let $A$ be a $2\times2$ matrix and $p_{A}(t)=t^{2}+5$. 

\begin{enumerate}
\item $t^{2}=-5$ has no real solutions. So $A$ has no real eigenvalues.
(it has the complex eigenvalues, $\pm i\sqrt{5}$)\Xspace\addtocounter{probs}{1}
\end{enumerate}
\end{enumerate}
\item (answers may vary) Give a definition for either \emph{eigenvalue }\uline{or}
a \emph{diagonalizable matrix (or linear transformation)}

\begin{itemize}
\item $\lambda$ is an \emph{eigenvalue} of $A\in\mathcal{M}_{n\times n}$
if and only if $A\mathbf{x}=\lambda\mathbf{x}$ for some nonzero vector
$\mathbf{x}\in\mathbb{R}^{\mbox{n}}$.\addtocounter{probs}{1}
\item $A$ is a \emph{diagonalizable} matrix if $A$ is similar to a diagonal
matrix, i.e., there exists a diagonal matrix $\Lambda$ such that
$A=P^{-1}\Lambda P$.\Xspace
\end{itemize}
\item {[}15 points{]} Let $A=\begin{bmatrix}\begin{array}{rr}
1 & 1\\
0 & -1
\end{array}\end{bmatrix}$. Find the eigenvalues and eigenvectors for $A$. 

\begin{enumerate}[resume]
\item []$p_{A}(t)=(1-t)(-1-t)$. So $A$ has eigenvalues $\lambda_{1}=1$
and $\lambda_{2}=-1$.
\item []$\mbox{E}(1):$ $\begin{bmatrix}\begin{array}{rr}
1-(1) & 1\\
0 & -1-(1)
\end{array}\end{bmatrix}=\begin{bmatrix}\begin{array}{rr}
0 & 1\\
0 & -2
\end{array}\end{bmatrix}\rightsquigarrow\begin{bmatrix}\begin{array}{rr}
0 & 1\\
0 & 0
\end{array}\end{bmatrix}\Rightarrow x_{2}=0$ and $x_{1}$ is arbitrary. So $\mbox{E}(1)$ has basis $\mathbf{v}_{1}=\begin{bmatrix}\begin{array}{r}
1\\
0
\end{array}\end{bmatrix}$
\item []$\mbox{E}(-1)$: $\begin{bmatrix}\begin{array}{rr}
1-(-1) & 1\\
0 & -1-(-1)
\end{array}\end{bmatrix}=\begin{bmatrix}\begin{array}{rr}
2 & 1\\
0 & 0
\end{array}\end{bmatrix}\Rightarrow2x_{1}=-x_{2}$ and so $\mbox{E}(-1)$ has basis $\mathbf{v}_{1}=\begin{bmatrix}\begin{array}{r}
1\\
-2
\end{array}\end{bmatrix}$\vfill\addtocounter{probs}{1}
\end{enumerate}
\item Prove the following statement: If $A$ and $B$ are similar matrices
then $\det(A)=\det(B)$. 

\begin{enumerate}[resume]
\item $A=P^{-1}BP$ since $A$ and $B$ are similar. So then $\det A=\det(P^{-1}BP)=\det P^{-1}\det B\det P=\det(P^{-1}P)\det B=\det(I_{n})\det B=\det B$.\vfill\addtocounter{probs}{1}
\end{enumerate}
\item {[}Bonus{]} Prove that if $A'$ is obtained from the $n\times n$
matrix $A$ by exchanging two rows then \\
$\det A'=-\det A$ (hint: use cofactors and induction). 

\begin{proof}
Let $n=2$. $\det\begin{bmatrix}\begin{array}{rr}
a & b\\
c & d
\end{array}\end{bmatrix}=ab-cd$ and $\det\begin{bmatrix}\begin{array}{rr}
c & d\\
a & b
\end{array}\end{bmatrix}=cb-ad=-(ad-cd)$. So it is true for $n=2$.\\
Now assume it is true for some $n\geq2$, and let $A\in\mathcal{M}_{(n+1)\times(n+1)}$
and $A'$ be obtained from $A$ by a row switch. $\det A'={\displaystyle \sum_{j=1}^{n+1}a_{ij}\det M_{ij}(A')}$
where $i$ is a row that was not switched ($n>2$ so such a row exists).
So then a rwo was switched in each $M_{ij}$ and so $\det M_{ij}(A')=-\det M_{ij}(A)$
for each minor matrix. Thus $\det A'={\displaystyle \sum_{j=1}^{n+1}-a_{ij}\det M_{ij}(A)=-\det A}$.\vfill
\end{proof}
\end{enumerate}

\end{document}
